%% ===========================================================================
\FloatBarrier
\section{Coefficients of the rectangular, elliptical and ellipsoidal averaging
windows}
\label{app:Mcoeff}
\setcounter{table}{0}
\setcounter{figure}{0}
%% ===========================================================================

The second moments collected here are the variances of the window coordinates
obtained by applying the average \eqref{eq:average} to the quadratic term of the
Taylor expansion \eqref{eq:expand}; higher-order window moments are excluded by this truncation. For a uniform kernel over a window $\mathcal{W}$ centered at
the origin, the normalized second moment, with $|\mathcal W|$ denoting area or volume, is
\begin{equation}
  M_{mn} = \frac{1}{|\mathcal{W}|}\int_{\mathcal W} s_m s_n \,\mathrm{d}\bm{s},
  \label{eq:appM}
\end{equation}
which is diagonal for a window whose axes are aligned with the coordinate axes
and which transforms as $M\mapsto RMR^{\mathsf T}$ under a rotation $R$ of the
window. The closed forms follow from
\eqref{eq:appM} by elementary integration, and the equivalence of windows with
the same $M$ is a statement about this truncated model: two kernels with the
same second moment can differ in higher moments, which are not included in the energy of \eqref{eq:energy}.

\paragraph{Cuboid} For $\mathcal W=[-a_1,a_1]\times[-a_2,a_2]$ the area is
$4a_1a_2$ and
\begin{equation}
  M_{11} = \frac{1}{4a_1a_2}\int_{-a_1}^{a_1}\!\!\int_{-a_2}^{a_2}
           s_1^{2}\,\mathrm{d}s_2\,\mathrm{d}s_1
         = \frac{1}{4a_1a_2}\cdot 2a_2\cdot\frac{2a_1^{3}}{3}
         = \frac{a_1^{2}}{3},
  \label{eq:Mcub}
\end{equation}
and identically $M_{22}=a_2^{2}/3$; the off-diagonal integral vanishes by the
oddness of $s_1s_2$. A three-dimensional cuboid of half-widths $a_i$ gives the
same coefficient along each axis.

\paragraph{Ellipse} For $\mathcal W=\{\bm{s}: (s_1/a_1)^{2}+(s_2/a_2)^{2}\le1\}$ the
area is $\pi a_1a_2$ and the polar second moment of the semi-axis is
$\int_{\mathcal W}s_1^{2}\,\mathrm{d}A = \pi a_1^{3}a_2/4$, so that
\begin{equation}
  M_{11} = \frac{\pi a_1^{3}a_2/4}{\pi a_1a_2} = \frac{a_1^{2}}{4},
  \qquad M_{22}=\frac{a_2^{2}}{4}.
  \label{eq:Mell}
\end{equation}

\paragraph{Ellipsoid} For the three-dimensional ellipsoid of semi-axes $a_i$
the volume is $4\pi a_1a_2a_3/3$ and
$\int_{\mathcal W}s_1^{2}\,\mathrm{d}V = 4\pi a_1^{3}a_2a_3/15$, hence
\begin{equation}
  M_{11} = \frac{a_1^{2}}{5},\qquad
  M_{22}=\frac{a_2^{2}}{5},\qquad
  M_{33}=\frac{a_3^{2}}{5}.
  \label{eq:Mellip}
\end{equation}

\paragraph{Equal-second-moment windows} Equations~\eqref{eq:Mcub}--\eqref{eq:Mellip}
are sufficient to explain shape equivalence at fixed second moment. A cuboid of
half-width $a$ and an ellipsoid of semi-axis $l$ in each compared direction carry the same second moment
when $a^{2}/3=l^{2}/5$, that is when $a=l\sqrt{3/5}$; for an ellipse the same
condition reads $a=l\sqrt{3/4}$. Two windows related in this way produce the
same gradient stiffness $K_g$ and therefore the same solution. The
window-shape and orientation calculations confirm the equivalence numerically,
while the ratio $5/3$ quantifies the difference for windows of equal half-width.

\paragraph{Rotated window} A rectangular window whose long axis makes an angle
$\theta$ with $x$ has
\begin{equation}
  \begin{aligned}
  M_{11}(\theta) &= M_{11}\cos^{2}\theta + M_{22}\sin^{2}\theta,\\
  M_{22}(\theta) &= M_{11}\sin^{2}\theta + M_{22}\cos^{2}\theta,\\
  M_{12}(\theta) &= \tfrac12\left(M_{11}-M_{22}\right)\sin2\theta.
  \end{aligned}
  \label{eq:Mrot}
\end{equation}
so that the off-diagonal coupling is largest at $\theta=45^\circ$ and vanishes
when the window is aligned with the axes. This is the tensor rotated in the orientation sweep.

\paragraph{Reduction to a scalar length} Let $L_{mn}$ denote the characteristic-length tensor of the one-scale
theory and $\ell$ its isotropic length. For the ellipsoid the corresponding coefficient is
$M_{mn}=L_{mn}/5$, hence $\tfrac14M_{mn}=\tfrac1{20}L_{mn}$ and the gradient
energy of~\cite{fem1} is recovered exactly, and setting
$M_{mn}=\ell^{2}\delta_{mn}$ reduces the same expression to the isotropic
Aifantis form, for which the boundary-layer length of a clamped bar is
$\sqrt{\ell^{2}/2}$. In the one-scale scalar notation
$L_{mn}=L\delta_{mn}$, $L$ denotes a squared characteristic length, not
the strip length; substituting $L=5\ell^{2}$ in $\sqrt{L/10}$ gives the
same boundary-layer length.

\FloatBarrier
